\section*{Chapter \ref{ch:iv1:speed-drills} --- Speed Drills}
\markboth{\protect\omTocMark{}{Answers: Speed Drills}}{\protect\omTocMark{}{Answers: Speed Drills}}

\begin{solution}{iq:iv1:speed-drills:1}
\textbf{Sheet 1.} Mark the sheet with this key. Your score is the number right minus twice the number wrong.
\par\nopagebreak
% Generated by code/interviews-1/03-speed-drills/python/gen_drill.py. Do not edit.
\begin{ivdrillkey}
\ivkey{1}{$1881$} & \ivkey{17}{$390$} & \ivkey{33}{$0.4$} & \ivkey{49}{$4.63$} & \ivkey{65}{$2760$} \\
\ivkey{2}{$750$} & \ivkey{18}{$75$} & \ivkey{34}{$0.27$} & \ivkey{50}{$2000$} & \ivkey{66}{$36$} \\
\ivkey{3}{$1071$} & \ivkey{19}{$20$} & \ivkey{35}{$3288$} & \ivkey{51}{$348$} & \ivkey{67}{$1675$} \\
\ivkey{4}{$570$} & \ivkey{20}{$1065$} & \ivkey{36}{$2590$} & \ivkey{52}{$\frac{5}{6}$}\rule[-9.5pt]{0pt}{23pt} & \ivkey{68}{$112$} \\
\ivkey{5}{$420$} & \ivkey{21}{$984$} & \ivkey{37}{$2499$} & \ivkey{53}{$1290$} & \ivkey{69}{$\frac{5}{3}$}\rule[-9.5pt]{0pt}{23pt} \\
\ivkey{6}{$26$} & \ivkey{22}{$\frac{13}{10}$}\rule[-9.5pt]{0pt}{23pt} & \ivkey{38}{$25$} & \ivkey{54}{$1241$} & \ivkey{70}{$195$} \\
\ivkey{7}{$23.86$} & \ivkey{23}{$43$} & \ivkey{39}{$20.96$} & \ivkey{55}{$3798$} & \ivkey{71}{$3168$} \\
\ivkey{8}{$2110$} & \ivkey{24}{$\frac{5}{12}$}\rule[-9.5pt]{0pt}{23pt} & \ivkey{40}{$12.4$} & \ivkey{56}{$224$} & \ivkey{72}{$893$} \\
\ivkey{9}{$55$} & \ivkey{25}{$\frac{2}{3}$}\rule[-9.5pt]{0pt}{23pt} & \ivkey{41}{$954$} & \ivkey{57}{$16$} & \ivkey{73}{$72$} \\
\ivkey{10}{$16$} & \ivkey{26}{$\frac{1}{8}$}\rule[-9.5pt]{0pt}{23pt} & \ivkey{42}{$5.8$} & \ivkey{58}{$540$} & \ivkey{74}{$1048$} \\
\ivkey{11}{$24.78$} & \ivkey{27}{$34$} & \ivkey{43}{$\frac{3}{25}$}\rule[-9.5pt]{0pt}{23pt} & \ivkey{59}{$1950$} & \ivkey{75}{$1053$} \\
\ivkey{12}{$\frac{2}{5}$}\rule[-9.5pt]{0pt}{23pt} & \ivkey{28}{$120$} & \ivkey{44}{$28$} & \ivkey{60}{$990$} & \ivkey{76}{$1448$} \\
\ivkey{13}{$627$} & \ivkey{29}{$473$} & \ivkey{45}{$1$} & \ivkey{61}{$\frac{5}{48}$}\rule[-9.5pt]{0pt}{23pt} & \ivkey{77}{$135$} \\
\ivkey{14}{$2191$} & \ivkey{30}{$204$} & \ivkey{46}{$42$} & \ivkey{62}{$1241$} & \ivkey{78}{$144$} \\
\ivkey{15}{$1355$} & \ivkey{31}{$1878$} & \ivkey{47}{$19$} & \ivkey{63}{$1211$} & \ivkey{79}{$0.04$} \\
\ivkey{16}{$3750$} & \ivkey{32}{$120$} & \ivkey{48}{$1280$} & \ivkey{64}{$2496$} & \ivkey{80}{$19$} \\
\end{ivdrillkey}


\textbf{When to leave a blank.} A guess that is right with probability $p$ adds $p - 2(1 - p) = 3p - 2$ to your score on average. That is positive only when $p > \tfrac23$, so leave an item blank unless you are fairly sure.

\textbf{The pace.} Eight minutes for 80 items is 6 seconds each. Skip an item as soon as you see no quick route, because a slow item costs you two or three easy ones.

\textbf{Where you lost points.} Each item belongs to one family. Find the family of each item you missed or skipped, and practise that family alone.
% Generated by code/interviews-1/03-speed-drills/python/gen_drill.py. Do not edit.
\ivfamily{Adding and subtracting whole numbers}{3, 8, 15, 16, 20, 31, 35, 41, 54, 55, 62, 72, 74, 75}{Work from the left, hundreds first, and keep a running total.}
\ivfamily{Adding and subtracting decimals}{7, 11, 39, 40, 42, 49, 57, 79}{Line up the decimal point in your head, or scale up by 100 and scale back.}
\ivfamily{Times one digit}{14, 21, 36, 48, 51, 60, 63, 65, 73, 76}{Split the larger number, or round up and subtract when that is shorter: $389 \times 7 = 400 \times 7 - 11 \times 7 = 2723$.}
\ivfamily{Numbers with a shortcut}{1, 2, 5, 13, 29, 37, 50, 53, 59, 64, 67, 71}{$\times 25$: times 100, then divide by 4. $\times 15$: add half and times 10. $\times 11$: times 10 and add itself. $\times 99$: times 100 and subtract itself. $(n - d)(n + d) = n^2 - d^2$.}
\ivfamily{Exact division}{6, 10, 27, 38, 44, 47, 66, 68, 70, 80}{$\div 5$: double and divide by 10. $\div 25$: times 4 and divide by 100. $\div 12$ and $\div 15$: divide by 3 first.}
\ivfamily{Fractions}{12, 22, 24, 25, 26, 43, 45, 52, 61, 69}{Use the smallest common denominator, multiply across, cancel before you multiply. Answers are in lowest terms.}
\ivfamily{Percentages}{4, 9, 17, 19, 28, 30, 32, 46, 56, 58}{Start from 10\% and build up: $15\%$ is $10\% + 5\%$, and $12.5\%$ is one eighth. $p\%$ of $n$ equals $n\%$ of $p$.}
\ivfamily{Decimal products}{18, 23, 33, 34, 77, 78}{Multiply the digits and count the places. Treat $0.25$, $0.75$, $1.5$ and $2.5$ as fractions: $\tfrac14$, $\tfrac34$, $\tfrac32$, $\tfrac52$.}


\iqwrong{Doing the items strictly in order and refusing to skip. The penalty punishes careless answers, not blanks, so a calm 60 right and none wrong scores better than 75 right and 8 wrong: 60 against 59.}
\iqfollow{How do I get a fresh sheet?}{Run \texttt{gen\_drill.py} in this chapter's code folder with another seed. Its items come from the same families, and it writes a new key.}
\end{solution}

\begin{solution}{iq:iv1:speed-drills:2}
\textbf{Sheet 2.} Mark the sheet with this key. Your score is the number right minus the number wrong.
\par\nopagebreak
% Generated by code/interviews-1/03-speed-drills/python/gen_drill.py. Do not edit.
\begin{ivdrillkey}
\ivkey{1}{$\frac{44}{35}$}\rule[-9.5pt]{0pt}{23pt} & \ivkey{11}{$8879$} & \ivkey{21}{$4225$} & \ivkey{31}{$9991$} & \ivkey{41}{$0.4375$} \\
\ivkey{2}{$13803$} & \ivkey{12}{$2025$} & \ivkey{22}{$\frac{27}{16}$}\rule[-9.5pt]{0pt}{23pt} & \ivkey{32}{$9936$} & \ivkey{42}{$2991$} \\
\ivkey{3}{$27.2$} & \ivkey{13}{$1925$} & \ivkey{23}{$\frac{7}{12}$}\rule[-9.5pt]{0pt}{23pt} & \ivkey{33}{$\frac{1}{3}$}\rule[-9.5pt]{0pt}{23pt} & \ivkey{43}{$93.5$} \\
\ivkey{4}{$85.5$} & \ivkey{14}{$\frac{37}{36}$}\rule[-9.5pt]{0pt}{23pt} & \ivkey{24}{$1829$} & \ivkey{34}{$14.4$} & \ivkey{44}{$0.45$} \\
\ivkey{5}{$0.128$} & \ivkey{15}{$5088$} & \ivkey{25}{$27.5$} & \ivkey{35}{$5988$} & \ivkey{45}{$4180$} \\
\ivkey{6}{$4376$} & \ivkey{16}{$3105$} & \ivkey{26}{$0.175$} & \ivkey{36}{$4489$} & \ivkey{46}{$0.3125$} \\
\ivkey{7}{$195.5$} & \ivkey{17}{$4439$} & \ivkey{27}{$31.5$} & \ivkey{37}{$3772$} & \ivkey{47}{$4649$} \\
\ivkey{8}{$0.375$} & \ivkey{18}{$2070$} & \ivkey{28}{$9135$} & \ivkey{38}{$\frac{65}{36}$}\rule[-9.5pt]{0pt}{23pt} & \ivkey{48}{$4096$} \\
\ivkey{9}{$9975$} & \ivkey{19}{$2736$} & \ivkey{29}{$4730$} & \ivkey{39}{$1517$} & \ivkey{49}{$\frac{10}{3}$}\rule[-9.5pt]{0pt}{23pt} \\
\ivkey{10}{$0.5625$} & \ivkey{20}{$2704$} & \ivkey{30}{$7474$} & \ivkey{40}{$1225$} & \ivkey{50}{$2030$} \\
\end{ivdrillkey}


\textbf{When to leave a blank.} At $+1$ and $-1$, a guess that is right with probability $p$ adds $p - (1 - p) = 2p - 1$ to your score on average. That is positive when $p > \tfrac12$, so the penalty here tolerates more doubt than on Sheet 1: answer when you think you are more likely right than wrong.

\textbf{The pace.} Six minutes for 50 items is 7.2 seconds each, and many items take a good calculator-free hand longer than that. Expect to leave some blank. The sheet rewards the order in which you take the items as well as your speed, so first pick off the items you can see through at once.

\textbf{Where you lost points.} Find the family of each item you missed or skipped.
% Generated by code/interviews-1/03-speed-drills/python/gen_drill.py. Do not edit.
\ivfamily{Adding and subtracting large numbers}{2, 6, 11, 15, 17, 28, 30, 47}{Round one number and correct: $8142 - 3678 = 8142 - 3700 + 22 = 4464$.}
\ivfamily{Two digits times two digits}{13, 16, 18, 19, 24, 29, 37, 39, 45, 50}{Split one factor, as in $37 \times 46 = 37 \times 40 + 37 \times 6 = 1480 + 222 = 1702$. If the factors are close, use the difference of squares instead.}
\ivfamily{Squares}{12, 20, 21, 36, 40, 48}{$n^2 = (n - d)(n + d) + d^2$, with $d$ chosen to make one factor round: $48^2 = 50 \times 46 + 4 = 2304$. A number ending in 5: $65^2$ is $6 \times 7$ followed by 25.}
\ivfamily{Near a round number}{9, 31, 32, 35, 42}{$103 \times 97 = 100^2 - 3^2$, and $998 \times 6 = 6000 - 12$.}
\ivfamily{Percentages of any number}{3, 4, 7, 25, 27, 34, 43}{Take 1\% or 10\% and build up, or swap the two: $24\%$ of 350 equals $350\%$ of 24, which is 84.}
\ivfamily{Fractions}{1, 14, 22, 23, 33, 38, 49}{Use the smallest common denominator. To divide, multiply by the upside-down fraction, and cancel first.}
\ivfamily{Decimal answers}{5, 8, 10, 26, 41, 44, 46}{Know the eighths and sixteenths: $\tfrac18 = 0.125$ and $\tfrac1{16} = 0.0625$. To divide by 40, divide by 4 and then by 10.}


\iqwrong{Slowing down to be sure of every item, as you would on Sheet 1. Here a slip costs only one point. A sheet with 30 right and 2 wrong scores 28, and so does 40 right and 12 wrong, so a faster and sloppier sitting loses nothing at that rate of error. Pace is worth more on this sheet than on the first.}
\iqfollow{Can I print the key with another mix?}{Yes. Run \texttt{gen\_drill.py} with sheet number 2 and another seed. It writes new items and a new key from the same families.}
\end{solution}

\begin{solution}{iq:iv1:speed-drills:3}
\textbf{Sheet 3.} Mark the sheet with this key, which gives the letter and the value. Your score is twice the number right minus the number wrong.
\par\nopagebreak
% Generated by code/interviews-1/03-speed-drills/python/gen_drill.py. Do not edit.
\begin{ivdrillkey}
\ivkey{1}{$\mathrm{D}\;68$} & \ivkey{9}{$\mathrm{D}\;92$} & \ivkey{17}{$\mathrm{A}\;9919$} & \ivkey{25}{$\mathrm{B}\;80$} & \ivkey{33}{$\mathrm{A}\;57393$} \\
\ivkey{2}{$\mathrm{C}\;21522$} & \ivkey{10}{$\mathrm{B}\;140.3$} & \ivkey{18}{$\mathrm{A}\;72$} & \ivkey{26}{$\mathrm{C}\;9964$} & \ivkey{34}{$\mathrm{D}\;385$} \\
\ivkey{3}{$\mathrm{C}\;9.12$} & \ivkey{11}{$\mathrm{D}\;30.87$} & \ivkey{19}{$\mathrm{D}\;6889$} & \ivkey{27}{$\mathrm{A}\;20010$} & \ivkey{35}{$\mathrm{C}\;18.2$} \\
\ivkey{4}{$\mathrm{A}\;48.75$} & \ivkey{12}{$\mathrm{B}\;52.9$} & \ivkey{20}{$\mathrm{B}\;26790$} & \ivkey{28}{$\mathrm{C}\;180.9$} & \ivkey{36}{$\mathrm{A}\;97$} \\
\ivkey{5}{$\mathrm{B}\;455$} & \ivkey{13}{$\mathrm{C}\;8228$} & \ivkey{21}{$\mathrm{B}\;8464$} & \ivkey{29}{$\mathrm{D}\;31.36$} &  \\
\ivkey{6}{$\mathrm{B}\;50.7$} & \ivkey{14}{$\mathrm{C}\;12160$} & \ivkey{22}{$\mathrm{C}\;60$} & \ivkey{30}{$\mathrm{C}\;34.2$} &  \\
\ivkey{7}{$\mathrm{D}\;67$} & \ivkey{15}{$\mathrm{D}\;43169$} & \ivkey{23}{$\mathrm{A}\;82.8$} & \ivkey{31}{$\mathrm{D}\;35$} &  \\
\ivkey{8}{$\mathrm{C}\;26.88$} & \ivkey{16}{$\mathrm{D}\;2916$} & \ivkey{24}{$\mathrm{A}\;44400$} & \ivkey{32}{$\mathrm{A}\;210$} &  \\
\end{ivdrillkey}


\textbf{When to leave a blank.} At $+2$ and $-1$, a guess that is right with probability $p$ adds $2p - (1 - p) = 3p - 1$ on average, so it pays when $p > \tfrac13$. A blind guess among four options loses $\tfrac14$ of a point. Among three it is worth nothing. Among two it gains $\tfrac12$. So guess as soon as you have removed two options, and stop to compute only when you have not.

\textbf{The pace.} Six minutes for 36 items is 10 seconds each, which is too little for most of the sums in full. Work on the options instead. Every item has three wrong options: one nearby with the wrong last digit, one with the right last digit and a slip in the tens, and most often one that is far from the right size. A last-digit check removes the first, and rounding removes the far one, which leaves a choice of two at the cost of a few seconds. The sum of the digits, taken modulo 9, then often settles it.

\textbf{Where you lost points.} Find the family of each item you missed or skipped.
% Generated by code/interviews-1/03-speed-drills/python/gen_drill.py. Do not edit.
\ivfamily{Three digits times two digits}{2, 13, 14, 15, 20, 24, 27, 33}{Do not finish the sum. Check the last digit, then the size, then the sum of the digits. For $378 \times 46$, the last digit is that of $8 \times 6$, and the size is near $400 \times 45$.}
\ivfamily{Squares and near-squares}{16, 17, 19, 21, 26}{$67^2$ ends in 9, and is near $70 \times 64$. $104 \times 96$ is $10000 - 16$.}
\ivfamily{Exact division}{1, 7, 18, 25, 31, 36}{Multiply back in your head: the quotient times the divisor must end in the right digit. Its size comes from rounding the divisor.}
\ivfamily{Percentages}{6, 10, 12, 23, 28, 30}{$17\%$ of 640 is $10\% + 5\% + 2\%$, that is $64 + 32 + 12.8$.}
\ivfamily{Decimal products}{3, 4, 8, 11, 29, 35}{Count the decimal places first. The options then differ only in digits you can test: $3.4 \times 1.7$ has two places and ends in 8.}
\ivfamily{Fractions of a number}{5, 9, 22, 32, 34}{Divide by the bottom first, then multiply by the top. For $\tfrac78$, take one eighth away.}


\iqwrong{Computing each item in full and checking the options last. Twenty items worked in full and sixteen left blank score 40. Working through the options, 24 right and 12 wrong score 36, so at that error rate the two are close, but a reader who removes options well gets far more than 24 right.}
\iqfollow{What if the real test has no options?}{Then Sheet 3's checks become your way of catching mistakes: the last digit, the size and the digit sum all still apply. Run \texttt{gen\_drill.py} with sheet number 3 and another seed for a new sheet.}
\end{solution}

\begin{solution}{iq:iv1:speed-drills:4}
\textbf{Sheet 4.} Mark the sheet with this key. Your score is the number right minus the number wrong.
\par\nopagebreak
% Generated by code/interviews-1/03-speed-drills/python/gen_drill.py. Do not edit.
\begin{ivdrillkey}
\ivkey{1}{$240$} & \ivkey{12}{$\frac{7}{144}$}\rule[-9.5pt]{0pt}{23pt} & \ivkey{23}{$65$} & \ivkey{34}{$765$} & \ivkey{45}{$\frac{23}{125}$}\rule[-9.5pt]{0pt}{23pt} \\
\ivkey{2}{$77.5$} & \ivkey{13}{$2.5$} & \ivkey{24}{$92$} & \ivkey{35}{$40$} & \ivkey{46}{$63$} \\
\ivkey{3}{$92.5$} & \ivkey{14}{$\frac{9}{20}$}\rule[-9.5pt]{0pt}{23pt} & \ivkey{25}{$40$} & \ivkey{36}{$\frac{39}{40}$}\rule[-9.5pt]{0pt}{23pt} & \ivkey{47}{$78$} \\
\ivkey{4}{$612$} & \ivkey{15}{$85$} & \ivkey{26}{$\frac{7}{20}$}\rule[-9.5pt]{0pt}{23pt} & \ivkey{37}{$\frac{1}{25}$}\rule[-9.5pt]{0pt}{23pt} & \ivkey{48}{$\frac{1}{2}$}\rule[-9.5pt]{0pt}{23pt} \\
\ivkey{5}{$60$} & \ivkey{16}{$365.5$} & \ivkey{27}{$\frac{13}{125}$}\rule[-9.5pt]{0pt}{23pt} & \ivkey{38}{$8.4$} & \ivkey{49}{$91.2$} \\
\ivkey{6}{$665$} & \ivkey{17}{$178.5$} & \ivkey{28}{$\frac{1}{125}$}\rule[-9.5pt]{0pt}{23pt} & \ivkey{39}{$14$} & \ivkey{50}{$\frac{19}{24}$}\rule[-9.5pt]{0pt}{23pt} \\
\ivkey{7}{$\frac{6}{7}$}\rule[-9.5pt]{0pt}{23pt} & \ivkey{18}{$214.5$} & \ivkey{29}{$260$} & \ivkey{40}{$55$} & \ivkey{51}{$1050$} \\
\ivkey{8}{$552$} & \ivkey{19}{$15$} & \ivkey{30}{$\frac{29}{40}$}\rule[-9.5pt]{0pt}{23pt} & \ivkey{41}{$8$} & \ivkey{52}{$10$} \\
\ivkey{9}{$\frac{1}{6}$}\rule[-9.5pt]{0pt}{23pt} & \ivkey{20}{$95$} & \ivkey{31}{$99$} & \ivkey{42}{$\frac{13}{35}$}\rule[-9.5pt]{0pt}{23pt} & \ivkey{53}{$156$} \\
\ivkey{10}{$\frac{26}{21}$}\rule[-9.5pt]{0pt}{23pt} & \ivkey{21}{$5$} & \ivkey{32}{$360$} & \ivkey{43}{$714$} & \ivkey{54}{$\frac{9}{7}$}\rule[-9.5pt]{0pt}{23pt} \\
\ivkey{11}{$35$} & \ivkey{22}{$50$} & \ivkey{33}{$\frac{71}{125}$}\rule[-9.5pt]{0pt}{23pt} & \ivkey{44}{$\frac{1}{20}$}\rule[-9.5pt]{0pt}{23pt} &  \\
\end{ivdrillkey}


\textbf{When to leave a blank.} At $+1$ and $-1$, a guess that is right with probability $p$ adds $2p - 1$ on average, so it pays when $p > \tfrac12$.

\textbf{The pace.} Six minutes for 54 items is under 7 seconds each. This sheet is about moving between three ways to write the same number, so the fastest items are the ones you know by heart. Take those first, and return to the others.

\textbf{Where you lost points.} Find the family of each item you missed or skipped.
% Generated by code/interviews-1/03-speed-drills/python/gen_drill.py. Do not edit.
\ivfamily{Fractions as percentages}{2, 3, 13, 20, 21, 24, 39, 40, 41, 47}{Know the eighths and the sixteenths. For a denominator of 20, 25 or 50, scale it up to 100: $\tfrac{7}{25} = \tfrac{28}{100}$. For 40, take half of the value over 20.}
\ivfamily{Decimals as fractions}{26, 27, 28, 30, 33, 36, 37, 45}{Read the decimal over a power of 10, then cancel: $0.375$ is $\tfrac{375}{1000}$, which is $\tfrac38$. Most of them are multiples of $\tfrac18$, $\tfrac1{20}$, $\tfrac1{25}$ or $\tfrac1{40}$.}
\ivfamily{Percentage change}{1, 4, 6, 8, 29, 32, 34, 43, 51, 53}{Add or subtract the percentage of the base, or multiply at once: $15\%$ up is $\times 1.15$, and $20\%$ down is $\times 0.8$.}
\ivfamily{As a percentage of}{11, 15, 19, 22, 23, 25, 35, 52}{Cancel first, then scale to 100: $36$ of $120$ is $\tfrac{3}{10}$, or $30\%$.}
\ivfamily{Percentages of any number}{5, 16, 17, 18, 31, 38, 46, 49}{Take 1\% or 10\% and build up, or swap the two: $p\%$ of $n$ equals $n\%$ of $p$.}
\ivfamily{Fractions}{7, 9, 10, 12, 14, 42, 44, 48, 50, 54}{Use the smallest common denominator. To divide, multiply by the upside-down fraction, and cancel first.}


\iqwrong{Doing a long division for every fraction. The eighths, the twentieths and the fortieths recur all over this sheet, and each one you know saves a few seconds on every later sheet.}
\iqfollow{What should I learn by heart first?}{The eighths: $12.5$, $37.5$, $62.5$ and $87.5$. Then the sixteenths, which are half of those: $\tfrac{3}{16}$ is $18.75\%$.}
\end{solution}

\begin{solution}{iq:iv1:speed-drills:5}
\textbf{Sheet 5.} Mark the sheet with this key. Your score is the number right.
\par\nopagebreak
% Generated by code/interviews-1/03-speed-drills/python/gen_drill.py. Do not edit.
\begin{ivdrillkey}
\ivkey{1}{$90$} & \ivkey{19}{$6$} & \ivkey{37}{$84$} & \ivkey{55}{$1736$} & \ivkey{73}{$41$} \\
\ivkey{2}{$60$} & \ivkey{20}{$91$} & \ivkey{38}{$382$} & \ivkey{56}{$128$} & \ivkey{74}{$91$} \\
\ivkey{3}{$20$} & \ivkey{21}{$210$} & \ivkey{39}{$92$} & \ivkey{57}{$1050$} & \ivkey{75}{$42$} \\
\ivkey{4}{$52$} & \ivkey{22}{$20$} & \ivkey{40}{$487$} & \ivkey{58}{$136$} & \ivkey{76}{$40$} \\
\ivkey{5}{$48$} & \ivkey{23}{$144$} & \ivkey{41}{$831$} & \ivkey{59}{$171$} & \ivkey{77}{$460$} \\
\ivkey{6}{$119$} & \ivkey{24}{$876$} & \ivkey{42}{$456$} & \ivkey{60}{$366$} & \ivkey{78}{$68$} \\
\ivkey{7}{$153$} & \ivkey{25}{$520$} & \ivkey{43}{$481$} & \ivkey{61}{$64$} & \ivkey{79}{$144$} \\
\ivkey{8}{$8$} & \ivkey{26}{$112$} & \ivkey{44}{$363$} & \ivkey{62}{$356$} & \ivkey{80}{$84$} \\
\ivkey{9}{$1134$} & \ivkey{27}{$68$} & \ivkey{45}{$61$} & \ivkey{63}{$85$} & \ivkey{81}{$190$} \\
\ivkey{10}{$844$} & \ivkey{28}{$1003$} & \ivkey{46}{$104$} & \ivkey{64}{$650$} & \ivkey{82}{$98$} \\
\ivkey{11}{$672$} & \ivkey{29}{$56$} & \ivkey{47}{$90$} & \ivkey{65}{$942$} & \ivkey{83}{$750$} \\
\ivkey{12}{$552$} & \ivkey{30}{$29$} & \ivkey{48}{$72$} & \ivkey{66}{$95$} & \ivkey{84}{$1128$} \\
\ivkey{13}{$74$} & \ivkey{31}{$36$} & \ivkey{49}{$539$} & \ivkey{67}{$402$} & \ivkey{85}{$726$} \\
\ivkey{14}{$1368$} & \ivkey{32}{$42$} & \ivkey{50}{$827$} & \ivkey{68}{$870$} & \ivkey{86}{$504$} \\
\ivkey{15}{$333$} & \ivkey{33}{$108$} & \ivkey{51}{$35$} & \ivkey{69}{$476$} & \ivkey{87}{$532$} \\
\ivkey{16}{$450$} & \ivkey{34}{$215$} & \ivkey{52}{$478$} & \ivkey{70}{$656$} & \ivkey{88}{$64$} \\
\ivkey{17}{$847$} & \ivkey{35}{$57$} & \ivkey{53}{$427$} & \ivkey{71}{$433$} & \ivkey{89}{$370$} \\
\ivkey{18}{$80$} & \ivkey{36}{$560$} & \ivkey{54}{$152$} & \ivkey{72}{$399$} & \ivkey{90}{$348$} \\
\end{ivdrillkey}


\textbf{When to leave a blank.} With no penalty, a guess costs nothing, so never leave a blank. If the time is about to run out, write the nearest value you have.

\textbf{The pace.} Seven minutes for 90 items is under 5 seconds each. That is too fast for most people to finish, however plain the items are. The sheet tests whether your simple arithmetic is automatic, because any thought at all uses up the time.

\textbf{Where you lost points.} Find the family of each item you missed or skipped.
% Generated by code/interviews-1/03-speed-drills/python/gen_drill.py. Do not edit.
\ivfamily{Adding and subtracting}{10, 11, 17, 24, 25, 28, 34, 36, 38, 40, 41, 42, 44, 49, 50, 53, 64, 66, 68, 83}{Work from the left, and do not write the carries.}
\ivfamily{Times tables to 19}{1, 2, 4, 5, 6, 7, 18, 20, 23, 26, 27, 31, 32, 33, 35, 37, 46, 48, 56, 58, 59, 63, 79, 80, 82}{Split the teens: $17 \times 8 = 80 + 56$.}
\ivfamily{Times one digit}{9, 12, 14, 15, 16, 54, 55, 57, 65, 70, 72, 84, 85, 86, 89}{Split the larger number, hundreds first.}
\ivfamily{Doubling and halving}{43, 52, 60, 62, 67, 69, 71, 77, 87, 90}{Go digit by digit from the left. These should take no thought at all.}
\ivfamily{Exact division}{13, 30, 39, 45, 51, 73, 74, 75, 78, 88}{Think of the times table backwards, and split a long number into easy parts: $476 \div 7 = 420 \div 7 + 56 \div 7$.}
\ivfamily{Easy percentages}{3, 8, 19, 21, 22, 29, 47, 61, 76, 81}{$10\%$ moves the point, $20\%$ doubles that, $25\%$ is a quarter and $50\%$ is a half.}


\iqwrong{Checking your work. With no penalty, a check earns nothing and costs time. A sitting with 70 items all right scores 70, and one with 90 items and 12 wrong scores 78. Where a wrong answer costs $-1$ or worse for each $+1$, the sums go the other way (at $+1$ and $-2$, 70 against 54), so read the scoring rule before you start.}
\iqfollow{What if I finish with time left?}{Then this level is automatic for you, and that is worth knowing. Go back over the items you skipped, and move on to the harder sheets.}
\end{solution}

\enlargethispage{-3\baselineskip}%
\begin{solution}{iq:iv1:speed-drills:6}
\textbf{Sheet 6.} Mark the sheet with this key. Your score is three times the number right minus the number wrong.
\par\nopagebreak
% Generated by code/interviews-1/03-speed-drills/python/gen_drill.py. Do not edit.
\begin{ivdrillkey}
\ivkey{1}{$\mathrm{D}\;9.48$} & \ivkey{9}{$\mathrm{B}\;9936$} & \ivkey{17}{$\mathrm{D}\;35$} & \ivkey{25}{$\mathrm{C}\;20.79$} & \ivkey{33}{$\mathrm{C}\;140$} \\
\ivkey{2}{$\mathrm{B}\;240.9$} & \ivkey{10}{$\mathrm{B}\;7938$} & \ivkey{18}{$\mathrm{D}\;62.1$} & \ivkey{26}{$\mathrm{B}\;462$} & \ivkey{34}{$\mathrm{D}\;9951$} \\
\ivkey{3}{$\mathrm{C}\;69126$} & \ivkey{11}{$\mathrm{C}\;66$} & \ivkey{19}{$\mathrm{A}\;58.52$} & \ivkey{27}{$\mathrm{A}\;81.4$} & \ivkey{35}{$\mathrm{D}\;35$} \\
\ivkey{4}{$\mathrm{D}\;21.32$} & \ivkey{12}{$\mathrm{D}\;455$} & \ivkey{20}{$\mathrm{A}\;476$} & \ivkey{28}{$\mathrm{C}\;44.24$} & \ivkey{36}{$\mathrm{B}\;56$} \\
\ivkey{5}{$\mathrm{A}\;13320$} & \ivkey{13}{$\mathrm{A}\;56$} & \ivkey{21}{$\mathrm{D}\;33927$} & \ivkey{29}{$\mathrm{B}\;66$} &  \\
\ivkey{6}{$\mathrm{D}\;49771$} & \ivkey{14}{$\mathrm{B}\;36.1$} & \ivkey{22}{$\mathrm{D}\;7225$} & \ivkey{30}{$\mathrm{C}\;61$} &  \\
\ivkey{7}{$\mathrm{C}\;37$} & \ivkey{15}{$\mathrm{D}\;14534$} & \ivkey{23}{$\mathrm{B}\;170.2$} & \ivkey{31}{$\mathrm{D}\;30120$} &  \\
\ivkey{8}{$\mathrm{D}\;7.14$} & \ivkey{16}{$\mathrm{A}\;72$} & \ivkey{24}{$\mathrm{D}\;6724$} & \ivkey{32}{$\mathrm{D}\;3844$} &  \\
\end{ivdrillkey}


\textbf{When to leave a blank.} At $+3$ and $-1$, a guess that is right with probability $p$ adds $4p - 1$ on average, so it pays when $p > \tfrac14$. A blind guess among four options is worth exactly 0. Remove one option, and a guess among three is worth $\tfrac13$. Remove two, and a guess between two is worth 1. So never leave an item blank if you can remove even one option.

\textbf{The pace.} Five minutes for 36 items is about 8 seconds each. As on Sheet 3, each item has one wrong option nearby, with a wrong last digit, one with a slip in the tens, and most often one of the wrong size.

\textbf{Where you lost points.} Find the family of each item you missed or skipped.
% Generated by code/interviews-1/03-speed-drills/python/gen_drill.py. Do not edit.
\ivfamily{Three digits times two digits}{3, 5, 6, 10, 15, 21, 31}{Test the last digit first, then the size.}
\ivfamily{Squares and near-squares}{9, 22, 24, 32, 34}{Check the last digit, then round. $(100 + d)(100 - d)$ is $10000 - d^2$.}
\ivfamily{Long division}{7, 13, 16, 17, 29, 30, 35}{Multiply back: the quotient's last digit times the divisor's last digit must give the last digit of the number.}
\ivfamily{Percentages}{2, 14, 18, 23, 27, 36}{Build the percentage from 10\%, 5\% and 1\%, and stop as soon as only one option is left.}
\ivfamily{Decimal products}{1, 4, 8, 19, 25, 28}{Count the places, then test the last digit of the product.}
\ivfamily{Fractions of a number}{11, 12, 20, 26, 33}{Divide by the bottom first, then multiply by the top.}


\iqwrong{Leaving blanks out of caution, as if the penalty were heavy. Removing one option on each of the 36 items and then guessing scores 12 on average, against 0 for blanks. The rule here rewards any information at all.}
\iqfollow{Why is the blind guess worth exactly zero?}{Because the rule was built that way: with four options, $+3$ and $-1$ average out to zero. If a test sets its penalty like this, you lose nothing by guessing, and you gain from anything you know.}
\end{solution}

\begin{solution}{iq:iv1:speed-drills:7}
\textbf{Sheet 7.} Mark the sheet with this key. An answer is right if it is within 5\% of the value in the key. The key gives the true value to three significant figures, which is plenty for this purpose. Your score is the number right minus the number wrong.
\par\nopagebreak
% Generated by code/interviews-1/03-speed-drills/python/gen_drill.py. Do not edit.
\begin{ivdrillkey}
\ivkey{1}{$85.3$} & \ivkey{7}{$1.34$} & \ivkey{13}{$1.65$} & \ivkey{19}{$76.2$} & \ivkey{25}{$465$} \\
\ivkey{2}{$13.2$} & \ivkey{8}{$97.9$} & \ivkey{14}{$9.84$} & \ivkey{20}{$17.6$} & \ivkey{26}{$1720$} \\
\ivkey{3}{$26.5$} & \ivkey{9}{$4050$} & \ivkey{15}{$1.28$} & \ivkey{21}{$1420$} & \ivkey{27}{$3.91$} \\
\ivkey{4}{$909$} & \ivkey{10}{$855$} & \ivkey{16}{$2.65$} & \ivkey{22}{$8.95$} & \ivkey{28}{$3440$} \\
\ivkey{5}{$6.68$} & \ivkey{11}{$17.6$} & \ivkey{17}{$2560$} & \ivkey{23}{$6.21$} & \ivkey{29}{$60.2$} \\
\ivkey{6}{$1070$} & \ivkey{12}{$1.63$} & \ivkey{18}{$7120$} & \ivkey{24}{$1.11$} & \ivkey{30}{$1.47$} \\
\end{ivdrillkey}


\textbf{When to leave a blank.} At $+1$ and $-1$, a guess that is right with probability $p$ adds $2p - 1$ on average, so it pays when $p > \tfrac12$. A rough estimate that you trust to within 5\% clears that bar easily. If you have only the order of magnitude, leave it blank.

\textbf{The pace.} Five minutes for 30 items is 10 seconds each, and exact work would take most of that for a single product. Round to numbers you can multiply, and correct only when the rounding moved them the same way.

\textbf{Where you lost points.} Find the family of each item you missed or skipped.
% Generated by code/interviews-1/03-speed-drills/python/gen_drill.py. Do not edit.
\ivfamily{Products}{6, 10, 17, 18, 21, 28}{Round one factor up and the other down, by about the same percentage: $38.7 \times 41.2$ is close to $40 \times 40$.}
\ivfamily{Quotients}{2, 3, 5, 14, 20, 22}{Round the divisor to an easy number, then correct by the percentage you moved it.}
\ivfamily{Square roots}{1, 8, 11, 19, 29}{Write the number as a square times a small factor, and move the root by half of the percentage gap: $\sqrt{5300}$ is about $70 + \tfrac{400}{140}$.}
\ivfamily{Compounding}{7, 12, 15, 16, 30}{Take $1.05^{10} \approx e^{0.5}$, and subtract a little, because the log of $1 + r$ is below $r$.}
\ivfamily{Exponentials and logarithms}{13, 23, 24, 27}{Use $e^x \approx 1 + x + \tfrac{x^2}{2}$ for a small $x$. Build each logarithm from $\ln 2$ and $\ln 10$.}
\ivfamily{Percentages}{4, 9, 25, 26}{Round the number, not the percentage, then move by the same small factor.}


\iqwrong{Calculating exactly. The product $38.7 \times 41.2$ is $1594.44$, and $40 \times 40 = 1600$ takes a second and is only $0.35\%$ away. Rounding both factors up, however, adds the two errors: $40 \times 45$ for $38.7 \times 42.6$ is $9\%$ too large.}
\iqfollow{What should I round first?}{Round in opposite directions if you can, since those errors cancel. If you cannot, note the size of the error and take that percentage off the result.}
\end{solution}

\enlargethispage{-1\baselineskip}%
\begin{solution}{iq:iv1:speed-drills:8}
\textbf{Sheet 8.} Mark the sheet with this key. Your score is twice the number right minus the number wrong.
\par\nopagebreak
% Generated by code/interviews-1/03-speed-drills/python/gen_drill.py. Do not edit.
\begin{ivdrillkey}
\ivkey{1}{$25$} & \ivkey{9}{$17$} & \ivkey{17}{$13$} & \ivkey{25}{$8$} & \ivkey{33}{$19$} \\
\ivkey{2}{$17$} & \ivkey{10}{$6$} & \ivkey{18}{$9$} & \ivkey{26}{$4$} & \ivkey{34}{$3$} \\
\ivkey{3}{$14$} & \ivkey{11}{$10$} & \ivkey{19}{$5$} & \ivkey{27}{$16$} & \ivkey{35}{$90$} \\
\ivkey{4}{$14$} & \ivkey{12}{$7$} & \ivkey{20}{$11$} & \ivkey{28}{$25$} & \ivkey{36}{$6$} \\
\ivkey{5}{$105$} & \ivkey{13}{$88$} & \ivkey{21}{$9$} & \ivkey{29}{$18$} & \ivkey{37}{$14$} \\
\ivkey{6}{$4$} & \ivkey{14}{$10$} & \ivkey{22}{$5$} & \ivkey{30}{$13$} & \ivkey{38}{$8$} \\
\ivkey{7}{$18$} & \ivkey{15}{$21$} & \ivkey{23}{$14$} & \ivkey{31}{$22$} & \ivkey{39}{$45$} \\
\ivkey{8}{$8$} & \ivkey{16}{$5$} & \ivkey{24}{$12$} & \ivkey{32}{$36$} & \ivkey{40}{$75$} \\
\end{ivdrillkey}


\textbf{When to leave a blank.} At $+2$ and $-1$, a guess that is right with probability $p$ adds $3p - 1$ on average, so it pays when $p > \tfrac13$. Every answer here is a whole number, so a fraction means you have slipped. Work it again if it is quick; otherwise leave it blank, since writing it is a sure $-1$.

\textbf{The pace.} Five minutes for 40 items is $7.5$ seconds each. Do not write the steps. Undo the operations in your head, and put the value back into the equation only if you have time.

\textbf{Where you lost points.} Find the family of each item you missed or skipped.
% Generated by code/interviews-1/03-speed-drills/python/gen_drill.py. Do not edit.
\ivfamily{Two steps}{7, 9, 15, 23, 24, 26, 27, 28, 31, 38}{Undo the addition, then the multiplication, in your head.}
\ivfamily{Brackets}{8, 11, 19, 21, 22, 30, 33, 37}{Divide by the factor first, if it goes in, and leave the bracket alone.}
\ivfamily{Both sides}{3, 4, 6, 12, 14, 16, 17, 36}{Gather the $x$ terms on one side, and the numbers on the other.}
\ivfamily{With a fraction}{5, 10, 13, 25, 32, 35, 39, 40}{Clear the fraction by multiplying, or read $\tfrac{12}{x} = 1.5$ as $x = \tfrac{12}{1.5}$.}
\ivfamily{Powers}{1, 2, 18, 20, 29, 34}{Know the squares up to 25 and the powers of 2 and 3.}


\iqwrong{Multiplying out every bracket. For $5(x - 2) = 35$, divide by 5 first and get $x - 2 = 7$. Multiplying out gives $5x - 10 = 35$, which is one more step for the same result.}
\iqfollow{Is a check by substitution worth the time?}{Only if it is quick. For $7x - 9 = 3x + 15$, the value $x = 6$ gives 33 on both sides, and that takes about as long as solving it. On a timed sheet, check only the answers that look odd.}
\end{solution}

\begin{solution}{iq:iv1:speed-drills:9}
\textbf{Sheet 9.} Mark the sheet with this key. Your score is twice the number right minus three times the number wrong.
\par\nopagebreak
% Generated by code/interviews-1/03-speed-drills/python/gen_drill.py. Do not edit.
\begin{ivdrillkey}
\ivkey{1}{$120$} & \ivkey{10}{$4096$} & \ivkey{19}{$625$} & \ivkey{28}{$2401$} & \ivkey{37}{$5$} \\
\ivkey{2}{$79$} & \ivkey{11}{$324$} & \ivkey{20}{$5040$} & \ivkey{29}{$16384$} & \ivkey{38}{$512$} \\
\ivkey{3}{$69$} & \ivkey{12}{$2197$} & \ivkey{21}{$14641$} & \ivkey{30}{$8192$} & \ivkey{39}{$3375$} \\
\ivkey{4}{$289$} & \ivkey{13}{$42$} & \ivkey{22}{$441$} & \ivkey{31}{$13$} & \ivkey{40}{$90$} \\
\ivkey{5}{$10$} & \ivkey{14}{$4$} & \ivkey{23}{$2048$} & \ivkey{32}{$7$} & \ivkey{41}{$46$} \\
\ivkey{6}{$15625$} & \ivkey{15}{$4$} & \ivkey{24}{$66$} & \ivkey{33}{$24$} & \ivkey{42}{$43$} \\
\ivkey{7}{$6$} & \ivkey{16}{$72$} & \ivkey{25}{$6$} & \ivkey{34}{$2187$} & \ivkey{43}{$14$} \\
\ivkey{8}{$256$} & \ivkey{17}{$125$} & \ivkey{26}{$529$} & \ivkey{35}{$44$} & \ivkey{44}{$1331$} \\
\ivkey{9}{$361$} & \ivkey{18}{$362880$} & \ivkey{27}{$7$} & \ivkey{36}{$625$} & \ivkey{45}{$32768$} \\
\end{ivdrillkey}


\textbf{When to leave a blank.} At $+2$ and $-3$, a guess that is right with probability $p$ adds $5p - 3$ on average, so it pays only when $p > \tfrac35$. Most of the items on this sheet are facts that you either know or do not. Leave a blank if you are not sure.

\textbf{The pace.} Five minutes for 45 items is under 7 seconds each. The sheet is a test of memory as much as of speed. Anything you must work out should go to a second pass.

\textbf{Where you lost points.} Find the family of each item you missed or skipped.
% Generated by code/interviews-1/03-speed-drills/python/gen_drill.py. Do not edit.
\ivfamily{Powers of two}{10, 23, 29, 30, 38, 45}{$2^{10} = 1024$. From there, double or halve, or read $2^{13}$ as $8 \times 1024$.}
\ivfamily{Other powers}{6, 12, 17, 19, 21, 28, 34, 39, 44}{Square, and multiply once more if you need to: $7^4 = 49^2$. $11^3 = 1331$ and $5^n$ is $10^n$ divided by $2^n$.}
\ivfamily{Squares up to 29}{4, 8, 9, 11, 22, 26, 36}{Use $n^2 = (n - d)(n + d) + d^2$, or simply know them.}
\ivfamily{Square roots}{2, 3, 13, 16, 24, 35, 41, 42}{The size gives the tens digit, and the last digit leaves at most two choices for the units. Compare with the square of the number ending in 5.}
\ivfamily{Cube roots}{14, 25, 31, 32, 43}{Know the cubes up to 15. The last digit of a cube fixes the last digit of its root.}
\ivfamily{Logarithms}{5, 7, 15, 27, 37}{Rewrite the number as a power of the base: $\log_2 4096$ is the power of two that makes 4096.}
\ivfamily{Factorials}{1, 18, 20, 33, 40}{Know them up to $9! = 362\,880$. A ratio such as $\tfrac{10!}{8!}$ is simply $10 \times 9$.}


\iqwrong{Building $2^{13}$ by doubling from 2. Start from $2^{10} = 1024$, and $2^{13}$ is $8 \times 1024 = 8192$. Knowing one landmark in each sequence does the work.}
\iqfollow{Which of these are worth learning by heart?}{The powers of 2 up to $2^{16} = 65\,536$, the squares up to 29, the cubes up to 15 and the factorials up to $9!$. They recur throughout the rest of this book.}
\end{solution}

\begin{solution}{iq:iv1:speed-drills:10}
\textbf{Sheet 10.} Mark the sheet with this key. Your score is the number right minus the number wrong.
\par\nopagebreak
% Generated by code/interviews-1/03-speed-drills/python/gen_drill.py. Do not edit.
\begin{ivdrillkey}
\ivkey{1}{$\frac{255}{256}$}\rule[-9.5pt]{0pt}{23pt} & \ivkey{9}{$120$} & \ivkey{17}{$55$} & \ivkey{25}{$\frac{1}{64}$}\rule[-9.5pt]{0pt}{23pt} & \ivkey{33}{$1287$} \\
\ivkey{2}{$90$} & \ivkey{10}{$56$} & \ivkey{18}{$62$} & \ivkey{26}{$66$} & \ivkey{34}{$\frac{5}{16}$}\rule[-9.5pt]{0pt}{23pt} \\
\ivkey{3}{$42$} & \ivkey{11}{$\frac{1}{512}$}\rule[-9.5pt]{0pt}{23pt} & \ivkey{19}{$33.5$} & \ivkey{27}{$\frac{63}{256}$}\rule[-9.5pt]{0pt}{23pt} & \ivkey{35}{$1$} \\
\ivkey{4}{$\frac{1023}{1024}$}\rule[-9.5pt]{0pt}{23pt} & \ivkey{12}{$23.5$} & \ivkey{20}{$336$} & \ivkey{28}{$18$} & \ivkey{36}{$84$} \\
\ivkey{5}{$\frac{2}{15}$}\rule[-9.5pt]{0pt}{23pt} & \ivkey{13}{$\frac{2}{7}$}\rule[-9.5pt]{0pt}{23pt} & \ivkey{21}{$30$} & \ivkey{29}{$39$} & \ivkey{37}{$\frac{15}{16}$}\rule[-9.5pt]{0pt}{23pt} \\
\ivkey{6}{$\frac{5}{32}$}\rule[-9.5pt]{0pt}{23pt} & \ivkey{14}{$78$} & \ivkey{22}{$182$} & \ivkey{30}{$\frac{1}{32}$}\rule[-9.5pt]{0pt}{23pt} & \ivkey{38}{$35$} \\
\ivkey{7}{$6$} & \ivkey{15}{$45$} & \ivkey{23}{$\frac{1}{32}$}\rule[-9.5pt]{0pt}{23pt} & \ivkey{31}{$\frac{21}{55}$}\rule[-9.5pt]{0pt}{23pt} & \ivkey{39}{$21$} \\
\ivkey{8}{$\frac{1}{26}$}\rule[-9.5pt]{0pt}{23pt} & \ivkey{16}{$\frac{6}{11}$}\rule[-9.5pt]{0pt}{23pt} & \ivkey{24}{$\frac{1}{21}$}\rule[-9.5pt]{0pt}{23pt} & \ivkey{32}{$\frac{5}{16}$}\rule[-9.5pt]{0pt}{23pt} & \ivkey{40}{$\frac{7}{64}$}\rule[-9.5pt]{0pt}{23pt} \\
\end{ivdrillkey}


\textbf{When to leave a blank.} At $+1$ and $-1$, a guess adds $2p - 1$ on average, so it pays when $p > \tfrac12$.

\textbf{The pace.} Six minutes for 40 items is 9 seconds each. This is the arithmetic that sits underneath a probability answer. The items are the kind you would otherwise slip on at the end of a correct argument.

\textbf{Where you lost points.} Find the family of each item you missed or skipped.
% Generated by code/interviews-1/03-speed-drills/python/gen_drill.py. Do not edit.
\ivfamily{Binomial coefficients}{9, 10, 14, 15, 17, 26, 33, 36, 38, 39}{Cancel before you multiply: $\binom{12}{3} = \tfrac{12 \cdot 11 \cdot 10}{6} = 2 \cdot 11 \cdot 10$.}
\ivfamily{Ratios of factorials}{2, 3, 20, 22}{Write only the factors that are left: $\tfrac{7!}{5!} = 7 \times 6$.}
\ivfamily{Powers of a half}{1, 4, 11, 23, 25, 37}{$2^{-5} = \tfrac{1}{32}$. The complement of $2^{-k}$ is $\tfrac{2^k - 1}{2^k}$.}
\ivfamily{Binomial chances}{6, 27, 30, 32, 34, 40}{Cancel the powers of two from the coefficient first, and give the result in lowest terms.}
\ivfamily{Drawing in a row}{5, 8, 13, 16, 24, 31}{Cancel across before you multiply: $\tfrac59 \times \tfrac48 = \tfrac59 \times \tfrac12$.}
\ivfamily{Expected values}{7, 12, 18, 19, 21, 28, 29, 35}{Multiply each chance by its amount, and keep the decimal point in place by treating $0.3$ as 3 tenths.}


\iqwrong{Multiplying out before you cancel. For $\binom{12}{3}$, computing $12 \times 11 \times 10 = 1320$ and dividing by 6 works, but $2 \times 11 \times 10 = 220$ needs no large number at all.}
\iqfollow{How can I check a binomial coefficient quickly?}{Use the step $\binom{n}{k} = \binom{n - 1}{k - 1} \cdot \tfrac{n}{k}$. For example, $\binom{10}{3} = \binom{9}{2} \cdot \tfrac{10}{3} = 36 \cdot \tfrac{10}{3} = 120$.}
\end{solution}

\begin{solution}{iq:iv1:speed-drills:11}
\textbf{Sheet 11.} Mark the sheet with this key. Your score is the number right minus three times the number wrong.
\par\nopagebreak
% Generated by code/interviews-1/03-speed-drills/python/gen_drill.py. Do not edit.
\begin{ivdrillkey}
\ivkey{1}{$91.2$} & \ivkey{6}{$31096$} & \ivkey{11}{$25149$} & \ivkey{16}{$356409$} & \ivkey{21}{$146.52$} \\
\ivkey{2}{$38624$} & \ivkey{7}{$93704$} & \ivkey{12}{$148$} & \ivkey{17}{$66.96$} & \ivkey{22}{$185.25$} \\
\ivkey{3}{$813604$} & \ivkey{8}{$23643$} & \ivkey{13}{$31710$} & \ivkey{18}{$27417$} & \ivkey{23}{$21614$} \\
\ivkey{4}{$9604$} & \ivkey{9}{$334$} & \ivkey{14}{$40804$} & \ivkey{19}{$29406$} & \ivkey{24}{$33222$} \\
\ivkey{5}{$315$} & \ivkey{10}{$51894$} & \ivkey{15}{$219$} & \ivkey{20}{$48654$} &  \\
\end{ivdrillkey}


\textbf{When to leave a blank.} At $+1$ and $-3$, a guess that is right with probability $p$ adds $4p - 3$ on average, so it pays only when $p > \tfrac34$. If you are not fairly sure of an item, leave it blank.

\textbf{The pace.} Six minutes for 24 items is 15 seconds each. This is an accuracy sheet. It is not meant to be finished at the expense of care, and working on paper is allowed.

\textbf{Where you lost points.} Find the family of each item you missed or skipped.
% Generated by code/interviews-1/03-speed-drills/python/gen_drill.py. Do not edit.
\ivfamily{Three digits times two digits}{2, 8, 10, 11, 18, 19, 20, 24}{Split the two-digit factor, and add carefully. Then check the last digit.}
\ivfamily{Near the hundreds}{6, 7, 13, 23}{$104 \times 297 = 104 \times 300 - 104 \times 3$.}
\ivfamily{Squares near the hundreds}{3, 4, 14, 16}{$(100k + d)^2 = 10000k^2 + 200kd + d^2$, with $d$ small and perhaps negative.}
\ivfamily{Long division}{5, 9, 12, 15}{Take off round multiples of the divisor, and keep a tally of the quotient.}
\ivfamily{Decimal products}{1, 17, 21, 22}{Count the places, multiply as whole numbers, and check the size against a rounded sum.}


\iqwrong{Hurrying to finish. A sitting with 24 items tried and 4 wrong scores $20 - 12 = 8$, and one with 15 tried and none wrong scores 15. At this penalty, one slip cancels three correct items.}
\iqfollow{How do I check an item without doing it twice?}{Check the last digit, and then the digit sum modulo 9. For $378 \times 46 = 17\,388$, the digit sums give $0 \times 1 = 0$, and the answer's digits add up to 27, which is 0 modulo 9. A digit-sum check misses only a swap of digits or an error of a multiple of 9.}
\end{solution}

\begin{solution}{iq:iv1:speed-drills:12}
\textbf{Sheet 12.} Mark the sheet with this key. Your score is the number right minus twice the number wrong, as on Sheet 1, so you can compare the two scores directly.
\par\nopagebreak
% Generated by code/interviews-1/03-speed-drills/python/gen_drill.py. Do not edit.
\begin{ivdrillkey}
\ivkey{1}{$4451$} & \ivkey{17}{$9996$} & \ivkey{33}{$4.8$} & \ivkey{49}{$570$} & \ivkey{65}{$14292$} \\
\ivkey{2}{$13$} & \ivkey{18}{$1575$} & \ivkey{34}{$792$} & \ivkey{50}{$2142$} & \ivkey{66}{$3552$} \\
\ivkey{3}{$65$} & \ivkey{19}{$1024$} & \ivkey{35}{$3484$} & \ivkey{51}{$\frac{32}{27}$}\rule[-9.5pt]{0pt}{23pt} & \ivkey{67}{$\frac{31}{50}$}\rule[-9.5pt]{0pt}{23pt} \\
\ivkey{4}{$289$} & \ivkey{20}{$12$} & \ivkey{36}{$6561$} & \ivkey{52}{$4225$} & \ivkey{68}{$6489$} \\
\ivkey{5}{$5$} & \ivkey{21}{$2409$} & \ivkey{37}{$2368$} & \ivkey{53}{$236.5$} & \ivkey{69}{$4$} \\
\ivkey{6}{$38$} & \ivkey{22}{$266.5$} & \ivkey{38}{$1588$} & \ivkey{54}{$125$} & \ivkey{70}{$21$} \\
\ivkey{7}{$1849$} & \ivkey{23}{$9975$} & \ivkey{39}{$7569$} & \ivkey{55}{$256$} & \ivkey{71}{$0.104$} \\
\ivkey{8}{$1148$} & \ivkey{24}{$2730$} & \ivkey{40}{$256$} & \ivkey{56}{$168$} & \ivkey{72}{$2398$} \\
\ivkey{9}{$\frac{7}{8}$}\rule[-9.5pt]{0pt}{23pt} & \ivkey{25}{$3988$} & \ivkey{41}{$99$} & \ivkey{57}{$35$} & \ivkey{73}{$4884$} \\
\ivkey{10}{$0.6875$} & \ivkey{26}{$1320$} & \ivkey{42}{$168$} & \ivkey{58}{$1408$} & \ivkey{74}{$9566$} \\
\ivkey{11}{$0.3125$} & \ivkey{27}{$56$} & \ivkey{43}{$495$} & \ivkey{59}{$12081$} & \ivkey{75}{$286$} \\
\ivkey{12}{$0.275$} & \ivkey{28}{$9984$} & \ivkey{44}{$8723$} & \ivkey{60}{$4990$} & \ivkey{76}{$0.4375$} \\
\ivkey{13}{$392$} & \ivkey{29}{$2952$} & \ivkey{45}{$11394$} & \ivkey{61}{$86$} & \ivkey{77}{$84$} \\
\ivkey{14}{$2916$} & \ivkey{30}{$4200$} & \ivkey{46}{$1850$} & \ivkey{62}{$8$} & \ivkey{78}{$70$} \\
\ivkey{15}{$10.2$} & \ivkey{31}{$16384$} & \ivkey{47}{$2209$} & \ivkey{63}{$3025$} & \ivkey{79}{$0.35$} \\
\ivkey{16}{$3432$} & \ivkey{32}{$19$} & \ivkey{48}{$330$} & \ivkey{64}{$\frac{1}{125}$}\rule[-9.5pt]{0pt}{23pt} & \ivkey{80}{$0.625$} \\
\end{ivdrillkey}


\textbf{When to leave a blank.} At $+1$ and $-2$, a guess that is right with probability $p$ adds $3p - 2$ on average, so it pays when $p > \tfrac23$.

\textbf{The pace.} Eight minutes for 80 items is 6 seconds each, the same as Sheet 1, with harder items. Few readers will finish it, so choose your items. Take each family you are strong in as soon as you see it, and skip the rest on the first pass.

\textbf{Where you lost points.} Find the family of each item you missed or skipped.
% Generated by code/interviews-1/03-speed-drills/python/gen_drill.py. Do not edit.
\ivfamily{Adding and subtracting large numbers}{1, 24, 37, 38, 44, 45, 59, 65, 68, 72, 74}{Round one number and correct.}
\ivfamily{Two digits times two digits}{8, 16, 18, 21, 26, 29, 30, 35, 46, 50, 58, 66, 73}{Split one factor, or use the difference of squares if the factors are close.}
\ivfamily{Squares}{7, 14, 39, 47, 52, 63}{$n^2 = (n - d)(n + d) + d^2$.}
\ivfamily{Near a round number}{17, 23, 25, 28, 60}{$(100 + d)(100 - d) = 10000 - d^2$.}
\ivfamily{Percentages}{4, 13, 15, 22, 27, 33, 40, 41, 42, 49, 53, 56}{Build from $10\%$ and $1\%$. For a change, multiply by $1 + r$.}
\ivfamily{Fractions and decimals}{9, 10, 11, 12, 51, 71, 76, 79, 80}{Know the eighths, sixteenths and fortieths. Cancel before you multiply.}
\ivfamily{Conversions}{3, 61, 64, 67}{Scale a fraction to 100 for a percentage. Read a decimal over a power of 10 for a fraction.}
\ivfamily{Powers and roots}{6, 19, 20, 31, 36, 54, 55, 57, 62, 77}{Know the powers of 2 up to $2^{16}$ and the squares up to 29.}
\ivfamily{Binomial coefficients}{34, 43, 48, 70, 75, 78}{Cancel before you multiply.}
\ivfamily{Algebra}{2, 5, 32, 69}{Undo each step in your head.}


\iqwrong{Attempting items at random until the time is up. Fifty right with none wrong scores 50, and 65 right with 10 wrong scores 45. Skipping is a skill here as much as speed is.}
\iqfollow{What is a good score?}{This book gives no norm, since it has none it can back up. Your own record is the useful measure: note your score against the date, return to the families you missed, and take a sheet from \texttt{gen\_drill.py} with a new seed in a week.}
\end{solution}
